% Full image and limit contribution to A.2(b,c), by /root/algebra_c.
% AI-produced by GPT-6 Astra (OpenAI), ChatGPT Pro mode, 19 November 2026.
\subsection{The exact image: simultaneous weighted conformal-radius inequalities}
\label{C:A2b-image-section}

Continue with a fixed stable pointed conformal surface $(X,p_1,\ldots,p_n)$ and its specified tangent coordinate norms. Let $\mathcal D_X$ be the systems of pairwise disjoint open conformal discs $D_i\ni p_i$, each avoiding every other marked point. Their boundaries need not be disjoint. Write $r_i(D)$ for their conformal radii and define, for $a\in\mathbb R_{\geq0}^n$,
\[
 H_X(a)=\sup_{D\in\mathcal D_X}\sum_i a_i\log r_i(D),\qquad H_X(0)=0.
\]
This definition is intrinsic once the tangent norms are fixed. The next statement computes the support function by a unique quadratic differential, so the eventual image test does not require a search over arbitrary discs.

\begin{theorem}[Known extremal theorem, with the complexity-zero case proved]
\label{C:A2b-support}
For $a\ne0$, let $Q_a$ be the unique Jenkins--Strebel differential whose characteristic domains are exactly the punctured discs at the active punctures $\{p_i:a_i>0\}$, with circumferences $2\pi\sqrt{a_i}$; at the inactive marked points it has at most simple poles. In particular no additional characteristic annuli are permitted. Let $D_i^a$ be its characteristic pointed discs and $R_i(a)$ their conformal radii. Then
\[
 H_X(a)=\sum_{a_i>0}a_i\log R_i(a)<\infty.
\]
Among disjoint active pointed discs avoiding all the other marked points, equality in this weighted radius inequality holds exactly for the indicated Strebel discs. Their complement is the critical graph, including inactive marked points, and has empty interior. The function $H_X$ is convex, positively homogeneous, and continuous on the whole nonnegative orthant.
\end{theorem}
\begin{proof}
For $3g+n-3>0$, existence and uniqueness are Liu (2008), Theorem~3.4; the extremal inequality and its equality clause are Lemma~1.7, pp.~218--219~\cite{C_Liu2008}. Take no characteristic annuli and circumference $2\pi\sqrt{a_i}$ at each active puncture. Lemma~1.2 identifies reduced modulus with $(2\pi)^{-1}\log r$. Substitution in Lemma~1.7 and division by $2\pi$ give exactly the displayed weighted inequality. Liu permits an arbitrary distinguished subset of the punctures, so inactive marked points remain excluded from all candidate discs. The stated graph complement is the characteristic-domain decomposition of this differential. Shrink each active disc by a factor $t<1$ in its Riemann coordinate and insert small discs at the inactive points. Letting $t\uparrow1$ proves that the supremum defining $H_X$ over complete systems has the same value as the active-disc extremum.

Liu assumes $3g+n-3>0$, which excludes the stable case $(0,3)$. We supply the exact reduction. Normalize its points to $0,1,\infty$ and take the double covering
\[
 \pi:\mathbb P^1_w\longrightarrow\mathbb P^1_z,\qquad z=w^2.
\]
The upstairs surface is marked at $0,1,-1,\infty$, so Liu applies. For downstairs weights $(a_0,a_1,a_\infty)$ use upstairs weights
\[
 (b_0,b_1,b_{-1},b_\infty)=(4a_0,a_1,a_1,4a_\infty).
\]
The involution $w\mapsto-w$ preserves these data. Uniqueness makes the upstairs Strebel differential invariant, so it descends to a quadratic differential downstairs. At a branch point invariance forces even Laurent powers, so an upstairs simple pole cannot occur. A regular value or an even-order zero descends to at most a simple pole; a double pole has its coefficient divided by four. Hence the descended differential has exactly the desired active circumferences and no extra singularities. Its characteristic discs descend, giving existence. Conversely the pullback of any downstairs Strebel differential has the specified upstairs characteristic domains; upstairs uniqueness gives downstairs uniqueness.

The same covering transfers the extremal inequality with no loss. A pointed disc at $0$ lifts to a branched degree-two disc with radius $\sqrt{r_0}$, and likewise at $\infty$ in coordinates $1/z,1/w$. A disc at $1$ lifts to two unbranched discs at $\pm1$, each of radius $r_1/2$. For example, if $f(0)=0$, the lift at $0$ is the single-valued univalent map $u\mapsto\sqrt{f(u^2)}$ with the simple zero and chosen derivative. Disjointness is preserved by full preimages. The upstairs weighted sum is therefore
\[
 2\bigl(a_0\log r_0+a_1\log r_1+a_\infty\log r_\infty-a_1\log2\bigr).
\]
The same identity holds for the Strebel radii. The additive constant cancels, proving the downstairs inequality. If equality holds, the upstairs equality clause identifies every active lifted disc, and projection identifies the downstairs discs. Zero weights simply omit the corresponding terms and marked discs from this argument. This covers every stable type.

Convexity and homogeneity follow from the supremum definition. For continuity, let $c_i=H_X(e_i)$ and translate all candidate logarithmic vectors by $-c=( -c_i)$, so their coordinates are nonpositive. Their support function is $H'(a)=H_X(a)-a\cdot c$. It is finite and lower semicontinuous as a supremum of continuous linear functions. If $a^{(m)}\to a\ne0$ in the nonnegative orthant, then for each $\epsilon>0$, eventually $a_i^{(m)}\geq(1-\epsilon)a_i$ for every $i$. For every nonpositive candidate vector $y$, this gives
$a^{(m)}\cdot y\leq(1-\epsilon)a\cdot y$. Taking suprema, then $\limsup$, and then $\epsilon\downarrow0$, yields $\limsup H'(a^{(m)})\leq H'(a)$. Lower semicontinuity gives equality. Continuity at $0$ follows from $H'\leq0$ and from a lower bound furnished by any single fixed small-disc system. Thus $H_X$ is continuous everywhere.
\end{proof}

\begin{proposition}[Proved: direct integral formula for the support data]
\label{C:A2b-support-integral}
The preceding support value can be evaluated directly from $Q_a$. On an active characteristic disc choose the branch
$\sqrt{Q_a}\sim i\sqrt{a_i}\,dz_i/z_i$ near $p_i$, and choose a regular point $b_i$ of its boundary. Then
\[
 w_i(z)=\exp\!\left(-\frac{i}{\sqrt{a_i}}
                  \int_{b_i}^{z}\sqrt{Q_a}\right)
\]
is its normalized disc coordinate up to rotation, and
\[
 \log R_i(a)=\lim_{z\to p_i}
 \left[\log|z_i(z)|-
 \operatorname{Re}\!\left(-\frac{i}{\sqrt{a_i}}
                  \int_{b_i}^{z}\sqrt{Q_a}\right)\right].
\]
Thus the image inequalities use only the unique meromorphic quadratic differential with principal parts $-a_i\,dz_i^2/z_i^2$ whose characteristic domains are precisely the punctured discs around the active marks, followed by these integrals and limits. At inactive marks at most simple poles are allowed; there are no additional characteristic annuli. No elementary expression for its coefficients in every genus is asserted.
\end{proposition}
\begin{proof}
The square root is single-valued on the punctured characteristic disc and has residue $i\sqrt{a_i}$. Its period around the puncture is $-2\pi\sqrt{a_i}$, which becomes $2\pi i$ in the exponential; hence $w_i$ is single-valued. Horizontal boundary arcs make the real part of the exponent zero, so $|w_i|=1$ there. In a cylinder coordinate $\zeta$ one has $Q_a=-a_i(d\zeta/\zeta)^2$ and the chosen square root is $i\sqrt{a_i}\,d\zeta/\zeta$. The displayed expression is therefore $\zeta/\zeta(b_i)$, proving it is a disc coordinate with a simple zero at $p_i$. Since $w_i=e^{i\alpha}z_i/R_i+O(z_i^2)$, taking logarithms gives the limit formula. Changing the boundary basepoint multiplies $w_i$ by a constant of modulus one and leaves $R_i$ unchanged.
\end{proof}

For a nonempty active index set $I\subset\{1,\ldots,n\}$, interpret $a\in\mathbb R_{\geq0}^I$ as zero on the other indices, and define the explicit convex domain
\[
 \mathcal K_{X,I}
 =\left\{y\in\mathbb R^I:
       \sum_{i\in I}a_i y_i<H_X(a)
       \quad\text{for every }a\geq0,\ \sum_{i\in I}a_i=1\right\}.
\]
It is nonempty: subtracting a positive constant from every coordinate of any small-disc logarithmic vector gives a point of this set. It is convex by the defining inequalities. It is open because the positive continuous gap $H_X(a)-a\cdot y$ has a positive minimum on the compact weight simplex.

\begin{proposition}[Proved: uniform support gap]
\label{C:A2b-gap}
For the canonical cap radii $F_X(x)_i=\log r_i$, with active parameters $x_i=\lambda_i^2>0$, we have
\[
 0\leq H_X(x)-x\cdot F_X(x)\leq2g-2+n.
\]
Here inactive ends remain cusps and are included in $n$. The lower inequality is strict for every nonzero finite parameter vector.
\end{proposition}
\begin{proof}
The nonnegative lower bound is Theorem~\ref{C:A2b-support}. For the upper bound, let $g$ be the Thurston metric of the canonical cap system, with hyperbolic core $M$ and boundary lengths $\ell_i=2\pi\lambda_i$. Gauss--Bonnet gives
\[
 A:=\operatorname{Area}_g(M)=2\pi(2g-2+n).
\]
The retraction that sends each flat cylinder $(u,y)$ to $(u,0)$ and is the identity on $M$ is length-nonincreasing. Every loop in the peripheral class of $p_i$ therefore has $g$-length at least $\ell_i$: after retraction it is a loop in the corresponding boundary class of the hyperbolic core. In its hyperbolic double, lifting such a loop joins $z$ to $\gamma_i z$, where $\gamma_i$ has translation length $\ell_i$, so its length is at least $\ell_i$. The fold retraction of the double to the core also shows that the peripheral class has not changed.

Use the Strebel metric $|Q_x|$, whose active cylinders have the same circumferences $\ell_i$. On one of these cylinders write its coordinates as
$(u,h)\in(\mathbb R/\ell_i\mathbb Z)\times(0,\infty)$ and write $g=\rho^2(du^2+dh^2)$. Every horizontal circle is peripheral, so
\[
 \int_0^{\ell_i}\rho(u,h)\,du\geq\ell_i.
\]
Cauchy--Schwarz gives $\int_0^{\ell_i}\rho^2\,du\geq\ell_i$. Integrating in $h$ and summing the cylinders compares the two areas on any union of complete truncated Strebel cylinders.

We state the finite-part calculation to make the cutoff comparison precise. Write $R_i=R_i(x)$ for the Strebel radii. Truncate its $i$th cylinder at
$T_i(\epsilon)=\lambda_i\log(R_i/\epsilon)$. Its flat area is
$\ell_iT_i=2\pi x_i\log(R_i/\epsilon)$. In the fixed local coordinate $z_i$, the removed end is bounded by
$|z_i|=\epsilon+O(\epsilon^2)$, since its normalized Riemann coordinate is $z_i/R_i+O(z_i^2)$. The area of the $O(\epsilon^2)$ discrepancy in either cylindrical metric is $O(\epsilon)$, hence tends to zero. In the canonical metric the corresponding finite part is
\[
 A+2\pi\sum_{i\in I}x_i\log(r_i/\epsilon)+o(1).
\]
Inactive cusps contribute finite area already counted in $A$. The Strebel critical graph has area zero for both metrics; its zeros and the inactive marked points cause no extra area term. Subtracting the common divergent terms from the area inequality and letting $\epsilon\to0$ gives
\[
 A+2\pi\bigl(x\cdot F_X(x)-H_X(x)\bigr)\geq0,
\]
which is the claimed upper bound. Strictness of the lower inequality follows either from the equality clause in Theorem~\ref{C:A2b-support}, since the canonical discs leave an open hyperbolic core, or from the explicit simultaneous collar enlargement proved next.
\end{proof}

\begin{proposition}[Proved: simultaneous enlargement and escape at zero length]
\label{C:A2b-collar}
Put
\[
 Y(\lambda)=\arctan\!\left(\frac{1}{2\sinh(\pi\lambda)}\right),
 \qquad \eta_i=Y(\lambda_i)/\lambda_i.
\]
The canonical cap discs can be simultaneously enlarged to disjoint pointed discs whose logarithmic radii are $F_X(x)_i+\eta_i$. Consequently
\[
 H_X(a)-a\cdot F_X(x)\geq\sum_{i\in I}a_i\eta_i>0
 \quad(a\geq0,\ a\ne0),
\]
and
\[
 F_X(x)_i\leq H_X(e_i)-Y(\lambda_i)/\lambda_i
        \longrightarrow-\infty\quad\text{as }\lambda_i\to0.
\]
The last estimate is uniform in all other lengths.
\end{proposition}
\begin{proof}
Mondello's published Collar Lemma~A.1 (2011, p.~668) gives disjoint embedded half-collars of area $\ell_i/[2\sinh(\ell_i/2)]$~\cite{C_Mondello2011}. In Fermi coordinates their area is $\ell_i\sinh s_i$, so the supplied width is
$s_i=\operatorname{arsinh}(1/[2\sinh(\pi\lambda_i)])$. Its conformal height is
\[
 \int_0^{s_i}\frac{ds}{\cosh s}
 =\arctan(\sinh s_i)=Y(\lambda_i).
\]
The canonical coordinate $w_i$ extends univalently from the flat cap through this collar to $|w_i|<\exp(Y(\lambda_i)/\lambda_i)$. The extended image is a pointed disc and its radius is multiplied by this exponential. The collar lemma gives disjointness and ensures that no other marked point is included. Applying the support inequality to this enlarged system proves the first estimate, and taking $a=e_i$ gives the second. Finally $Y(\lambda)\to\pi/2$ as $\lambda\downarrow0$, proving uniform escape to $-\infty$.
\end{proof}

\begin{theorem}[Proved: complete image and uniqueness, A.2(b)]
\label{C:A2b-image}
For every stable fixed pointed conformal surface and every nonempty active set $I$, the map
\[
 F_X:(0,\infty)^I\longrightarrow\mathcal K_{X,I}
\]
is a real-analytic diffeomorphism. Thus a vector of positive radii $(r_i)_{i\in I}$ is realized by a unique admissible canonical cap system, with all other ends cuspidal, if and only if
\[
 \prod_{i\in I}r_i^{a_i}
 <\prod_{a_i>0}R_i(a)^{a_i}
 \quad\text{for every }a\geq0,\quad\sum_{i\in I}a_i=1.
\]
Here $R_i(a)$ are computed from the unique Jenkins--Strebel system of circumferences $2\pi\sqrt{a_i}$. The zero radius vector corresponds to the all-cuspidal system. These statements give the whole image on all faces, and injectivity on their union.
\end{theorem}
\begin{proof}
The collar estimate puts $F_X(x)$ in $\mathcal K_{X,I}$. The energy theorem and Corollary~\ref{C:A2b-injectivity} make it injective and a local real-analytic diffeomorphism, so its image is open. We prove that the image is closed relative to $\mathcal K_{X,I}$.

Suppose $F_X(x^{(m)})\to y\in\mathcal K_{X,I}$. Compactness of the weight simplex and continuity of $H_X$ give
\[
 \delta:=\min_{a\geq0,\ \sum a_i=1}\bigl(H_X(a)-a\cdot y\bigr)>0.
\]
For all large $m$, the same gap with $y$ replaced by $F_X(x^{(m)})$ is at least $\delta/2$, uniformly in $a$. Homogeneity and Proposition~\ref{C:A2b-gap} then imply
\[
 (\delta/2)\sum_i x_i^{(m)}
 \leq H_X(x^{(m)})-x^{(m)}\cdot F_X(x^{(m)})
 \leq2g-2+n.
\]
The parameters are bounded above. No active coordinate can tend to zero, because Proposition~\ref{C:A2b-collar} would force its logarithmic radius to $-\infty$, contradicting convergence to $y$. Therefore a subsequence converges to $x$ in the positive orthant, and continuity gives $F_X(x)=y$.

The nonempty image is now both open and closed in the connected convex set $\mathcal K_{X,I}$, so it is all of it. Injectivity glues the local analytic inverses into a global inverse. Exponentiating the defining inequalities gives exactly the displayed intrinsic radius test. Different faces are distinguished by their zero entries, and the all-cuspidal face is a point. No unproved convexity assertion about arbitrary disc systems was used: the image convexity in logarithmic coordinates follows from this argument.
\end{proof}

\subsection{The full large-length limit, A.2(c)}

\begin{lemma}[Proved: compactness of normalized embedded discs]
\label{C:A2b-normal-families}
For a fixed index $i$, let $f_m:\mathbb D\to X$ be injective holomorphic maps with $f_m(0)=p_i$ whose images avoid every $p_j$, $j\ne i$. Normalize their derivatives to be positive in the chosen coordinate norm. Their conformal radii have a finite upper bound. Every sequence has a subsequence converging locally uniformly to a holomorphic map; if its radii are bounded away from zero, the limit is an injective holomorphic disc avoiding the other marked points. Nonconstant limits of mutually disjoint discs are mutually disjoint. If the radii converge to a positive value, convergence of the normalized maps to a nonconstant limit is equivalent to pointed conformal kernel convergence of their image domains.
\end{lemma}
\begin{proof}
We use Koebe's quarter and growth estimates in the form proved in Siu (2005), Appendix, pp.~10--13~\cite{C_Siu2005}. We spell out normality in all genera. If $g=0$, stability gives at least two marked points other than $p_i$. Take a fixed global coordinate sending $p_i$ to $0$ and two of the other points to $1,\infty$. The maps become univalent plane maps omitting $1$ and fixing $0$. Koebe's quarter theorem gives $|f_m'(0)|\leq4$. Its distortion bound
\[
 |f_m(w)|\leq |f_m'(0)|\frac{|w|}{(1-|w|)^2}
\]
gives local uniform boundedness. The derivative in the original coordinate differs by one fixed nonzero factor.

If $g=1$, choose a flat uniformization $\pi:\mathbb C\to X$ and normalize the lifts $\tilde f_m$ by $\tilde f_m(0)=0$. They are univalent. Since $f_m$ is an embedding and $\pi$ is locally isometric for the flat metric,
\[
 \int_{\mathbb D}|\tilde f_m'|^2\,dA
 =\operatorname{Area}_{\mathrm{flat}}(f_m(\mathbb D))
 \leq\operatorname{Area}_{\mathrm{flat}}(X).
\]
The mean-value inequality bounds $|\tilde f_m'(0)|$, and the same univalent distortion theorem gives local boundedness. If $g\geq2$, lift instead to the universal covering disc, with a fixed lift of $p_i$ at the origin; the lifted maps are bounded by one and form a normal family. In each case the normalized derivatives identify the limiting radius, and a positive limiting radius excludes a constant limit.

For completeness, the univalent-limit and disjointness assertions do not require any boundary regularity. If a nonconstant limit identified two distinct parameter points, take disjoint small parameter discs around them and a coordinate near their common image. Rouch\'e's theorem would force the nearby injective maps to take the same value on both discs, a contradiction. The same local argument excludes a critical point, an omitted marked point in the limiting image, or an intersection of two nonconstant limiting discs.

Let $D=f(\mathbb D)$ for a nonconstant normalized limit. Rouch\'e's theorem on small parameter discs shows that every compact subset of $D$ is eventually contained in $D_m=f_m(\mathbb D)$. Conversely, suppose $\Omega$ is a connected pointed open domain every compact subset of which is eventually contained in $D_m$. The inverse maps $f_m^{-1}$ are bounded holomorphic functions on those compact subsets. A diagonal normal-family argument gives $h:\Omega\to\overline{\mathbb D}$. Its derivative at $p_i$ is the positive reciprocal of the limiting radius, so it is nonconstant and $h(\Omega)\subset\mathbb D$. Passing to the limit in $f_m\circ f_m^{-1}=\operatorname{id}$ yields $f\circ h=\operatorname{id}_\Omega$. Hence $\Omega\subset D$, which proves the maximality required for kernel convergence. The reverse implication follows by the same normal-family argument and uniqueness of the normalized Riemann map.
\end{proof}

\begin{theorem}[Proved from the support-gap estimate: the complete projective limit]
\label{C:A2b-strebel-limit}
Let $x^{(m)}=(\lambda_i^{(m)})^2$ be a sequence of squared length vectors with
\[
 s_m=\sum_i x_i^{(m)}\longrightarrow\infty,
 \qquad a^{(m)}=x^{(m)}/s_m\longrightarrow a,
 \qquad a_i\geq0,\quad\sum_i a_i=1.
\]
The canonical cap at every index with $a_i>0$ converges, in the pointed conformal-kernel sense, to its unique Jenkins--Strebel disc with circumference proportional to $\sqrt{a_i}$. Every remaining cap has conformal radius tending to zero and has degenerate pointed kernel. In particular, if one boundary length tends to infinity while the others remain bounded, its cap tends to the unique one-weight Strebel slit domain. This is a domain limit on the fixed conformal surface; it does not assert Hausdorff convergence of arbitrary boundary traces.
\end{theorem}
\begin{proof}
Write $y_i^{(m)}=\log r_i^{(m)}$ for positive-length caps. Pass to a subsequence on which the set $J$ of positive-length indices is fixed; there are only finitely many such sets. Omit the indices outside $J$ throughout. Every index with $a_i>0$ belongs to $J$. Proving the asserted limit for every subsequence with fixed zero pattern proves it for the original sequence. Proposition~\ref{C:A2b-gap} gives
\[
 0\leq H_X(x^{(m)})-\sum_i x_i^{(m)}y_i^{(m)}\leq C,
 \qquad C=2g-2+n.
\]
Homogeneity gives
\begin{equation}
 \sum_i a_i^{(m)}y_i^{(m)}\geq H_X(a^{(m)})-C/s_m.
 \label{C:eq-radius-support-limit}
\end{equation}
Every coordinate has the fixed upper bound
$y_i^{(m)}\leq M_i:=H_X(e_i)<\infty$. A single fixed system of small disjoint discs also bounds $H_X$ below uniformly on the weight simplex. If $a_i>0$, then $a_i^{(m)}$ is bounded away from zero; isolating its term in~\eqref{C:eq-radius-support-limit} and using the other upper bounds gives a uniform lower bound on $y_i^{(m)}$.

Put $I=\{i:a_i>0\}$. Lemma~\ref{C:A2b-normal-families} therefore gives, along a subsequence, mutually disjoint nonconstant limiting discs $D_i$ for all $i\in I$, and convergent logarithmic radii $y_i$. Dropping the inactive terms using their upper bounds in~\eqref{C:eq-radius-support-limit}, we obtain
\[
 \sum_{i\in I}a_i y_i
 \geq\liminf_m H_X(a^{(m)})\geq H_X(a).
\]
The last inequality requires only lower semicontinuity: $H_X$ is a supremum of continuous linear functions, so it is lower semicontinuous. The weighted extremal theorem applied to the limiting disjoint active discs gives the reverse inequality. Equality follows, and its uniqueness clause identifies every $D_i$ as the designated Jenkins--Strebel disc.

These active Strebel discs cover $X$ except for their critical graph and the inactive marked points, a set with empty interior. If some inactive cap had radius bounded below along a further subsequence, Lemma~\ref{C:A2b-normal-families} would give an additional nonconstant open disc disjoint from all the active limiting discs. This is impossible. Hence every inactive radius tends to zero. No fixed open neighborhood of its center can lie eventually in such cap domains, by monotonicity of conformal radius, so their pointed kernels are degenerate.

Finally every convergent subsequence of normalized active cap maps has the same limit by uniqueness of the weighted extremal system and the derivative normalization. Normal-family compactness therefore makes the whole sequence converge. Lemma~\ref{C:A2b-normal-families} converts this to the asserted pointed kernel convergence. When only the $i$th boundary length diverges and the others stay bounded, $a=e_i$, giving precisely the one-weight slit domain.
\end{proof}

\begin{proposition}[Proved: explicit pants limits and failure of the box model]
\label{C:A2b-pants-extremes}
On $X=\mathbb P^1$ with marks $0,1,\infty$, use coordinates $z,z-1,1/z$ respectively. The one-weight extremal disc at $0$ is $\mathbb C\setminus[1,\infty)$, with radius $4$. For equal weights at $0,1$ and zero weight at $\infty$, the extremal discs are the half-planes $\operatorname{Re}z<1/2$ and $\operatorname{Re}z>1/2$, both of radius $1$. Thus every finite canonical cap system satisfies $r_0r_1<1$, although the individual supremum of each of these radii is $4$. Its image is therefore not a box. The limiting disc at $0$ is the slit domain when $\lambda_0\to\infty$ with the others bounded; it is the indicated half-plane when $\lambda_0=\lambda_1\to\infty$ and $\lambda_\infty$ stays bounded.
\end{proposition}
\begin{proof}
For one active weight take
\[
 Q=\frac{dz^2}{z^2(z-1)},\qquad
 z=\frac{4u}{(1+u)^2}.
\]
Direct substitution gives $Q=-du^2/u^2$ and the map has derivative $4$ at $0$. It maps the unit disc univalently onto $\mathbb C\setminus[1,\infty)$: its inverse is $(1-\sqrt{1-z})/(1+\sqrt{1-z})$, where the square root has positive real part off the slit, and the boundary formula $4e^{it}/(1+e^{it})^2=1/\cos^2(t/2)$ identifies its slit boundary. This supplies the one-weight characteristic disc and its radius.

For the two equal active weights, take
\[
 Q=-\frac{dz^2}{z^2(z-1)^2},\qquad u=\frac{z}{1-z}.
\]
Then $Q=-du^2/u^2$. The regions $|u|<1$ and $|u|>1$ correspond exactly to $\operatorname{Re}z<1/2$ and $\operatorname{Re}z>1/2$. Their normalized disc maps are $u\mapsto u/(1+u)$ and $v\mapsto1/(1+v)$, with derivative moduli $1$. The marked point at infinity lies on their separating trajectory. The support theorem gives $H_X(1,1,0)=0$, so the strict image inequality is $r_0r_1<1$. The one-weight computation gives $H_X(e_0)=\log4$; symmetry gives the same individual supremum at $1$. Shrinking either slit disc and inserting small discs at the other two marks shows that every radius below $4$ occurs as the corresponding coordinate of some interior admissible radius vector. Thus the simultaneous restriction cannot be a product of the individual intervals. The two rational pullback identities and radius constants are independently checked in the attached exact symbolic program. The two limits follow from Theorem~\ref{C:A2b-strebel-limit}, whose limiting squared-length weights are respectively $(1,0,0)$ and $(1/2,1/2,0)$.
\end{proof}

\paragraph{Original two-domain source.}
The same non-strict radius product bound for arbitrary disjoint simply
connected planar domains is Lavrentiev's original Theorem~3
\citep[\S2, pp.~169--173]{Lavrentiev1934}, now opened in full.
His maps $F_i:D_i\to\D$, $F_i(z_i)=0$, satisfy
$|F_0'(z_0)|\,|F_1'(z_1)|\ge |z_1-z_0|^{-2}$.
Since $r(D_i,z_i)=|F_i'(z_i)|^{-1}$, this gives
$r(D_0,z_0)r(D_1,z_1)\le|z_1-z_0|^2$.
For the canonical finite caps, our simultaneous collar enlargement
proves the strict inequality. Thus the originally remembered
Lavrentiev input is verified directly and its equality case is
distinguished from the interior admissible image.
